% ED26-REQ-N-c-004-02: exact STOC 1982 extended-abstract presentation.

\subsection{Vardi 1982: query complexity and ordered fixed-point capture}
\label{sec:p09-n-c-004-02}
\index{data complexity}\index{expression complexity}\index{combined complexity}

Vardi's extended abstract compares the cost of answering a fixed query with
the cost of evaluating an input expression.  Those are different families of
recognition problems.  Its fixed-point capture statement is a qualified
remark using the source's \(YF^+\), not a definition of the later general
syntax FO(LFP) \cite[pp.~137--143]{Vardi1982}.

\begin{hbdefinition}{def:p09-vardi82-query}{Database, query and recognition graph}
A database is \(B=(D,R_1,\ldots,R_k)\), where \(D\subseteq\mathbb N\)
is finite and \(R_i\subseteq D^{a_i}\) with \(a_i>0\).
Its type is the fixed rank vector \(\bar a=(a_1,\ldots,a_k)\).
A query of type \(\bar a\to b\) is a partial function on databases of
that type whose defined value at \(B\) is a subset of \(D^b\).
The computable queries in the source satisfy the consistency criterion:
renaming the domain transports both the input relations and the query answer.
A language \(L\) is a set of expressions with a meaning function
\(\mu\); write \(Q_e=\mu(e)\) and
\(C(L)=\{Q_e:e\in L\}\).

The recognition graph of a query and the expression graph of a fixed database
are respectively
\[
 \operatorname{Gr}(Q)=\{(\bar d,B):\bar d\in Q(B)\},\qquad
 \operatorname{Gr}_{L}(B)=
    \{(\bar d,e):e\in L,\ \bar d\in Q_e(B)\}.
\]
Membership includes that the query is defined.  This is tuple recognition,
not a claim that all answer tuples are supplied in constant time.
For a class of languages considered here, the expression fixes the answer
rank, so enumerating the at most \(\lvert D\rvert^b\) candidates accounts
for answer production on the data axis.  With expression rank varying,
that enumeration is not a fixed-degree polynomial bound.
Boolean answers can be read as membership of the
empty tuple in \(D^0\).  The source assumes a standard explicit encoding
of the finite domain and relation tables, and gives a tuple-list encoding
example; no succinct circuit representation is assumed
\cite[\S2, pp.~138--139]{Vardi1982}.
\end{hbdefinition}

\begin{hbdefinition}{def:p09-vardi82-axes}{The three input axes and completeness}
\emph{Data complexity} studies \(\operatorname{Gr}(Q_e)\) for each fixed
\(e\), with database and candidate tuple varying.
\emph{Expression complexity} studies \(\operatorname{Gr}_L(B)\) for each
fixed \(B\), with expression and candidate tuple varying.
\emph{Combined complexity} allows all three to vary and measures the
combined encoded input size.  It is introduced in the paper, but the paper's
numbered classifications concentrate on the first two axes.

The language is D-complete in a class \(\mathcal K\) when every fixed
expression's query graph is in \(\mathcal K\) and one fixed expression
has a graph to which every language in \(\mathcal K\) logspace-reduces.
It is E-complete when every fixed database's expression graph is in
\(\mathcal K\) and one fixed database gives such a hard expression graph.
Finally
\(Q\mathcal K=\{Q:\operatorname{Gr}(Q)\in\mathcal K\}\)
is a class of queries, not a class of expressions
\cite[pp.~137, 139]{Vardi1982}.
For capture statements below write \(Q\mathcal K_{\rm tot}\) for its
subclass of total consistent queries.
\end{hbdefinition}

\begin{hbwarning}{def:p09-vardi82-graph-domain}{An answer-membership graph does not encode definedness}
The nowhere-defined partial query and the query defined everywhere with the
empty answer have the same empty recognition graph.  They are different
partial functions.  The source's logical expressions define total queries,
so an equality with a graph-complexity class of \emph{all partial functions}
would be false under literal function equality.  The capture statements here
therefore use total queries.  For a partial query with an efficiently
recognizable graph, the construction yields a total query with the same
answer-membership graph; it does not reconstruct the partial function's
domain of definition.
\end{hbwarning}

\begin{hbdefinition}{def:p09-vardi82-languages}{The source languages and the positive operator}
The source \(F\) consists of expressions \(\bar x.\varphi(\bar x)\)
where \(\varphi\) is first-order with equality and no functions, and the
distinct output variables \(\bar x\) include all its free variables.
It returns the tuples satisfying \(\varphi\).  Its quantifier-free part is
\(F^-\).  The extension \(TF\) adds
\(\tau.(x,y).\varphi(x,y)\), returning the positive-length transitive
closure of the binary relation given by the first-order expression.

The extension \(YF\) adds expressions
\(\Upsilon R.\bar x.\varphi(\bar x)\), with the rank of \(R\) equal
to the number of output slots and every occurrence of \(R\) positive in
the first-order \(\varphi\): after expanding abbreviations, it is under
an even number of negations.  Its value is the least fixed point of
\[
 \Phi_B(X)=\{\bar d\in D^{\lvert\bar x\rvert}:
                    (B,X)\models\varphi(\bar d)\}.
\]
The least fixed point is obtained from the empty relation by iteration.
These source definitions add the stated outer operators to first-order
expressions; they are not silently replaced here by an unrestricted nested
fixed-point grammar.  The source \(SF\) similarly adds existential
quantification of auxiliary relations to first-order expressions.
The source identifies the positive fixed-point queries as total
\cite[pp.~139--140]{Vardi1982}.

In Remark~2, \(YF^+\) additionally permits existential quantification over
conjunctions of equality conditions with \(YF\) expressions.  This provides
restriction and projection of a fixed-point relation.  For example, projecting
away tableau coordinates subject to equality of two flag slots is a use of
\(YF^+\), not of the bare expression \(\Upsilon R.\bar x.\varphi\).
In the concluding section the source distinguishes the logically closed
language \(Y\), obtained by closing \(YF\) under Boolean connectives and
first-order quantification, and reports \(C(Y)=C(YF^+)\), citing Immerman's
independent work in the same 1982 proceedings \cite[p.~145]{Vardi1982}.
This source-cited normal form does not make the initial grammar logically
closed by definition. In particular, when empty universes are admitted as
in the database definition above, it needs the qualification proved in
\cref{thm:p09-vardi82-empty-boundary}; it is not used here to bypass that
boundary.
\end{hbdefinition}

\begin{hbtheorem}{thm:p09-vardi82-theorem5}{proved=annex:p9.n-c-004.2.theorem5}{Original Theorem 5: two different complete problems}
\(YF\) is D-complete in \(PTIME\) and E-complete in \(EXPTIME\),
with the logspace reduction convention of
Definition~\ref{def:p09-vardi82-axes}.  For E-hardness the source fixes
\(B_0=(\{0,1\},\{\langle1\rangle\})\), with a unary relation naming
the element~1 \cite[Theorem~5, pp.~141--142]{Vardi1982}.
\end{hbtheorem}

Vardi uses Path System Accessibility for D-hardness. His E-hardness proof
describes a \((2n+m)\)-ary relation on \(B_0\): two blocks of
\(n\) Boolean slots encode time and tape position, and \(m\) more slots
encode the fixed machine's cell symbol.  The fixed point simulates an
exponential-length deterministic computation while the \emph{database stays
fixed}.  Increasing the relation arity is possible because the expression is
now input.  This is not the lower-bound argument for a fixed query on varying
databases.  The source explicitly says that its proofs in this version are
brief sketches \cite[pp.~138, 141--142]{Vardi1982}. The annex supplies
complete reductions: a polynomial-time machine is translated to a path
system by a logspace enumerator, while an exponential-time machine is
translated to a positive expression and one target tuple on the fixed
Boolean database.

\begin{hblemma}{thm:p09-vardi82-iteration}{proved=annex:p9.n-c-004.2.iteration}{The positive-iteration bound on all three axes}
For a fixed positive expression whose defined relation has rank \(r\),
iteration stabilizes after at most \(\lvert D\rvert^r\) strict increases,
followed by an equality check.  Its tuple-recognition data complexity is
polynomial time.  With expression size \(s\) varying and database fixed,
the straightforward evaluation is at most exponential in a polynomial in
\(s\); with both varying it is at most exponential in a polynomial in
their combined explicit input size.  These upper bounds do not identify data
complexity with expression or combined complexity.
\end{hblemma}

\begin{hbexample}{def:p09-vardi82-flights}{Flights with one company}
Let \(R_2(c,u,v)\) say that company \(c\) has a direct flight from
\(u\) to \(v\).  The source expression
\[
 \Upsilon R.(c,u,v).
   \bigl(R_2(c,u,v)\lor
          \exists w(R_2(c,u,w)\land R(c,w,v))\bigr)
\]
returns paths of positive length using one fixed company
\cite[p.~140]{Vardi1982}.  For direct flights
\(R_2(a,p,q)\), \(R_2(a,q,r)\), \(R_2(b,p,q)\), its first stage
contains these three triples and its second stage adds \((a,p,r)\).
It does not add \((b,p,r)\): a flight of company~a cannot complete a
path whose company slot is~b.  A path's company coordinate must be preserved
through the recursion.  The bare \(TF\) operator defined in this paper takes
an unparameterized binary first-order relation, which is why this ternary
parameter-preserving example motivates the larger source language.
\end{hbexample}

\begin{hbtheorem}{thm:p09-vardi82-remark2}{proved=annex:p9.n-c-004.2.ordered}{Remark 2: total-query capture with the domain convention explicit}
On nonempty finite databases with a distinguished input linear order
\(\leq\),
\[
 C(YF^+,\leq)=QPTIME_{\rm tot}(\leq).
\]
If empty universes are also allowed, the same construction covers every
positive answer rank, and every Boolean polynomial-time query that is
false on the empty database. For all total queries, including Boolean
queries true there, the closed language satisfies
\[
 C(Y,\leq)=QPTIME_{\rm tot}(\leq).
\]
On all finite databases, without a supplied order,
\[
 C(SF)=QNPTIME_{\rm tot}.
\]
These are equalities of total consistent
queries; alternatively they may be read at tuple-recognition-graph scope.
The order in the deterministic statements is supplied by the database.
The nondeterministic statement may guess an order as part of its
existential relation certificate
\cite[Remark~2, pp.~142--143]{Vardi1982}.
\end{hbtheorem}

The annex supplies a single positive deterministic tableau relation, including
tuple-valued queries, canonical relabeling, first-order addressing,
initialization, local updates, singleton inputs, and answer projection.
Its existential-second-order argument uses the proved original Fagin
theorem, explicitly converts tuple markers to free output variables,
combines the finite auxiliary relation family into one relation, and handles
singleton and empty inputs. These are reconstructed proofs of the source's
mechanisms. They do not claim that the extended abstract prints these
details or the later normal-form proof.

\begin{hblemma}{thm:p09-vardi82-empty-boundary}{proved=annex:p9.n-c-004.2.empty}{The empty-domain boundary of the literal projection grammar}
Every Boolean \(YF^+\) expression in
\cref{def:p09-vardi82-languages} which is true on an empty universe is
equivalent to a first-order sentence. Consequently, with empty universes
included, this literal grammar does not capture all ordered polynomial-time
Boolean queries. Even cardinality, counting zero as even, is a counterexample.
The closed language \(Y\) does express it by adjoining the empty case to
the nonempty capture construction.
\end{hblemma}

\begin{hbexample}{def:p09-vardi82-even}{EVEN separates unordered capture from data completeness}
Vardi's \(EVEN\) is the Boolean query true exactly when \(\lvert D\rvert\)
is even.  It is in \(QLOGSPACE\), but is not in \(C(YF^+)\) on
unordered databases; the source attributes nondefinability to Chandra and
Harel \cite[Remark~2, p.~143]{Vardi1982}.
The handbook's equality-pattern proof is
Theorem~\ref{thm:p09-fagin74-unordered-boundary}.  Every source \(YF^+\)
expression embeds in order-free FO(LFP) by replacing its positive operator
with a least-fixed-point subformula and retaining its equality conditions and
existential projections.  That stronger logic already fails on parity of pure
sets, or of structures whose input relations are all empty.  Thus D-completeness
in \(PTIME\) does not mean that every unordered polynomial-time query is
expressible.  An encoding reduction can place a hard problem into a query
without making the language capture all queries of that complexity class.
\end{hbexample}

\begin{table}[htbp]
\centering\small
\begin{tabular}{@{}p{.20\textwidth}p{.32\textwidth}p{.38\textwidth}@{}}
\toprule
Task & Fixed and varying inputs & Applicable source result or boundary \\
\midrule
Repeatedly run one recursive query & Expression fixed; database and candidate
tuple varying & Theorem~5 D-completeness in \(PTIME\); positive rank \(r\)
is fixed. \\
Evaluate many expressions on one database & Database fixed; expression and
candidate tuple varying & Theorem~5 E-completeness in \(EXPTIME\); increasing
arity can encode exponentially many tuple addresses even on \(B_0\). \\
Evaluate an arbitrary expression on an arbitrary database & Both vary; total
encoded size measured & This is combined complexity, not either of the preceding
families. The iteration lemma supplies an upper bound, not a mislabeled
Theorem~5 combined-completeness assertion. \\
Define a particular ordered polynomial-time query & Schema and query fixed;
ordered database varies & On nonempty domains Remark~2 uses \(YF^+\),
equality restrictions and projection. Closed \(Y\) also handles Boolean
queries true on the empty domain. \\
Define parity without order & Domain varies; order absent & \(EVEN\) remains
logspace computable but is outside unordered \(C(YF^+)\). \\
\bottomrule
\end{tabular}
\caption{Vardi's input-axis lookup: completeness and capture are different claims.}
\label{tab:p09-vardi82-input-axes}
\end{table}
